\section*{Chapter \ref{ch:b3:modules} --- Modules over a Principal
Ideal Domain}

\begin{solution}{exo:b3:modules:1}
(a) If $\Q = \Z q_1 + \dots + \Z q_k$, let $d$ be a common
denominator of the $q_i$: every combination lies in $\frac1d\Z$,
but $\frac1{2d} \notin \frac1d\Z$. Contradiction.

(b) Torsion-free: $nq = 0$ with $n \neq 0$ forces $q = 0$ in
$\Q$. Not free: any two nonzero rationals $\frac ab, \frac cd$
satisfy the nontrivial relation $(bc)\frac ab - (ad)\frac cd =
0$, so a basis has at most one element; $\Q \cong \Z$ would make
$\Q = \Z q$ cyclic, but $\frac q2 \notin \Z q$. (And $\Q \neq
0$.)

(c) \cref{thm:b3:modules:structure} assumes \emph{finite
generation}, which (a) denies: no contradiction --- rather, $\Q$
shows the hypothesis is necessary in the statement
``torsion-free $\Rightarrow$ free''.
\end{solution}

\begin{solution}{exo:b3:modules:2}
$360 = 2^3\cdot3^2\cdot5$. Partitions: of $3$: $(3), (2,1),
(1,1,1)$; of $2$: $(2), (1,1)$; of $1$: $(1)$. Hence $3 \times 2
\times 1 = 6$ groups. Elementary divisors $\to$ invariant
factors:
\begin{center}
\begin{tabular}{l l}
$\Z/8 \times \Z/9 \times \Z/5$ & $\Z/360$\\
$\Z/8 \times \Z/3 \times \Z/3 \times \Z/5$ & $\Z/3 \times
\Z/120$\\
$\Z/4 \times \Z/2 \times \Z/9 \times \Z/5$ & $\Z/2 \times
\Z/180$\\
$\Z/4 \times \Z/2 \times \Z/3 \times \Z/3 \times \Z/5$ & $\Z/6
\times \Z/60$\\
$(\Z/2)^3 \times \Z/9 \times \Z/5$ & $\Z/2 \times \Z/2 \times
\Z/90$\\
$(\Z/2)^3 \times \Z/3 \times \Z/3 \times \Z/5$ & $\Z/2 \times
\Z/6 \times \Z/30$
\end{tabular}
\end{center}
(To pass to invariant factors: the largest $d_s$ collects the
highest prime power of each prime, and so on down.) Of order
$p^5$: as many as partitions of $5$, namely $7$.
\end{solution}

\begin{solution}{exo:b3:modules:3}
$B$: $D_1 = \gcd(2,4,6,8) = 2$; $D_2 = \abs{\det B} = \abs{16 -
24} = 8$. Invariant factors $d_1 = 2$, $d_2 = 8/2 = 4$: Smith
form $\operatorname{diag}(2, 4)$, and $\Z^2/B\Z^2 \cong \Z/2\Z
\times \Z/4\Z$.

$C$: diagonal but \emph{not} Smith ($2 \nmid 3$). $D_1 =
\gcd(2,3,12) = 1$; $D_2 = \gcd(2\cdot3,\, 2\cdot12,\, 3\cdot12) =
\gcd(6, 24, 36) = 6$; $D_3 = 72$. So $d = (1, 6, 12)$ and
$\Z^3/C\Z^3 \cong \Z/6\Z\times\Z/12\Z$ --- consistently with the
CRT: $\Z/2\times\Z/3\times\Z/12 \cong \Z/6\times\Z/12$.
\end{solution}

\begin{solution}{exo:b3:modules:4}
Write $B = Q\,\operatorname{diag}(d_1, \dots, d_n)\,P$ with $Q, P
\in GL_n(\Z)$ (\cref{thm:b3:modules:smith}; no zero $d_i$ since
$\det B \neq 0$). Then $\Z^n/B\Z^n \cong \Z^n/
\operatorname{diag}(d)\Z^n = \prod_i \Z/d_i\Z$ (the composed
isomorphism $x \mapsto Q^{-1}x$ of $\Z^n$ maps $B\Z^n$ onto
$\operatorname{diag}(d)P\Z^n = \operatorname{diag}(d)\Z^n$). Its
cardinality is $\prod \abs{d_i} = \abs{\det
\operatorname{diag}(d)} = \abs{\det B}$, as $\det Q, \det P = \pm
1$. For $L = \Z(2,0) + \Z(1,3)$: $B = \bigl(\begin{smallmatrix}2
& 1\\ 0 & 3\end{smallmatrix}\bigr)$, $D_1 = 1$, $D_2 = 6$: $\Z^2/L
\cong \Z/6\Z$, of cardinality $\abs{\det B} = 6$.
\end{solution}

\begin{solution}{exo:b3:modules:5}
(a) Submodule: done in \cref{def:b3:modules:torsion}. If $a(x +
T(M)) = 0$ in $M/T(M)$ with $a \neq 0$, then $ax \in T(M)$: $bax
= 0$ for some $b \neq 0$, and $ba \neq 0$ (domain), so $x \in
T(M)$: the class is zero. $M/T(M)$ is torsion-free.

(b) $(X, Y) \subseteq K[X,Y]$ is torsion-free (a submodule of the
domain $K[X,Y]$ acting on itself). Suppose it were free; any two
elements $P, Q$ satisfy $Q\cdot P - P \cdot Q = 0$, a nontrivial
relation when $P \ne Q$ are nonzero, so a basis has one element:
$(X, Y) = (P)$ principal --- contradicting
\cref{exo:b3:rings:6}(a). Torsion-free and finitely generated
($X, Y$ generate), yet not free: over the non-PID $K[X,Y]$, the
structure theorem fails.
\end{solution}

\begin{solution}{exo:b3:modules:6}
(a) If $\Z = 2\Z \oplus C$, the projection $\Z \to \Z/2\Z$
restricts to an isomorphism $C \cong \Z/2\Z$: $C$ would be a
subgroup of $\Z$ whose nonzero element $x$ satisfies $2x \in C
\cap 2\Z = 0$. But $\Z$ is torsion-free: $C = 0$, forcing $\Z =
2\Z$ --- false.

(b) $A^n/M$ is finitely generated and torsion-free, hence free
(\cref{thm:b3:modules:structure}): $A^n/M \cong A^r$ with basis
$f_1, \dots, f_r$. Choose preimages $y_i \in A^n$ of the $f_i$
and set $F = Ay_1 + \dots + Ay_r$. Every $x \in A^n$ has $\pi(x)
= \sum a_if_i$, so $x - \sum a_iy_i \in M$: $A^n = M + F$. If
$\sum a_iy_i \in M$, applying $\pi$ gives $\sum a_if_i = 0$,
hence all $a_i = 0$ (basis): $M \cap F = 0$. So $A^n = M \oplus
F$.
\end{solution}

\begin{solution}{exo:b3:modules:7}
The reduction of \cref{exo:b3:modules:3} was effective: with
\[
L = \begin{pmatrix} 1 & 0\\ -3 & 1\end{pmatrix},
\qquad
R = \begin{pmatrix} 1 & 2\\ 0 & -1 \end{pmatrix},
\qquad
LBR = \begin{pmatrix} 2 & 0\\ 0 & 4\end{pmatrix}
\]
(row operation $R_2 \leftarrow R_2 - 3R_1$, column operations
$C_2 \leftarrow C_2 - 2C_1$ then $C_2 \leftarrow -C_2$). Setting
$y = R^{-1}x$ (a bijection of $(\Z/20\Z)^2$, $R$ being
invertible over $\Z$), the system $Bx \equiv b \pmod{20}$ is
equivalent to
\[
2y_1 \equiv b_1, \qquad 4y_2 \equiv -3b_1 + b_2 \pmod{20}.
\]
The congruence $ky \equiv c \pmod{20}$ is solvable iff
$\gcd(k, 20) \mid c$: solutions exist iff $2 \mid b_1$ and $4
\mid b_2 - 3b_1$, i.e.\ \emph{$b_1$ even and $b_2 \equiv 3b_1
\pmod 4$}. When solvable there are $2 \times 4 = 8$ solutions
modulo $20$.
\end{solution}

\begin{solution}{exo:b3:modules:8}
(a) A nilpotent $u$ has $\mu_u = X^k$; the elementary divisors
are $X^{k_1} \geq \dots$, one Jordan block $J_{k_i}(0)$ per part
of a partition of $4$:
\begin{center}
\begin{tabular}{l l l}
partition & Jordan form & invariant factors\\
$(4)$ & $J_4$ & $X^4$\\
$(3,1)$ & $J_3 \oplus J_1$ & $X,\ X^3$\\
$(2,2)$ & $J_2 \oplus J_2$ & $X^2,\ X^2$\\
$(2,1,1)$ & $J_2 \oplus J_1 \oplus J_1$ & $X,\ X,\ X^2$\\
$(1,1,1,1)$ & $0$ & $X, X, X, X$
\end{tabular}
\end{center}
(b) Take $u = J_2\oplus J_2$ and $v = J_2 \oplus J_1 \oplus J_1$:
both have $\chi = X^4$, $\mu = X^2$, but different invariant
factors --- not similar (\cref{thm:b3:modules:frobenius}); one
can also compare ranks: $\operatorname{rk} u = 2 \neq 1 =
\operatorname{rk} v$. For $n \leq 3$: $\chi$ and $\mu$ determine,
for each eigenvalue $\lambda$ (over a splitting field), the total
size $m_\lambda \leq 3$ of the $\lambda$-blocks and the largest
block $r_\lambda$; a partition of $m \leq 3$ is determined by its
largest part ($m = 3, r = 2$ forces $(2,1)$, etc.). So the
elementary divisors coincide, and
\cref{cor:b3:modules:descent} descends the similarity to the base
field.
\end{solution}

\begin{solution}{exo:b3:modules:9}
(i)$\Rightarrow$(ii): if $V = K[u]x$, then $V \cong
K[X]/\operatorname{Ann}(x)$, and $\operatorname{Ann}(x) = (\mu_u)$
(a polynomial kills $x$ iff it kills all of $V = K[u]x$, since
$P(u)Q(u)x = Q(u)P(u)x$). So $n = \dim V = \deg \mu_u$; as
$\mu_u \mid \chi_u$ and $\deg\chi_u = n$, monicity gives $\mu_u =
\chi_u$.

(ii)$\Rightarrow$(iii): $\deg\chi_u = \sum_i \deg P_i$ and
$\mu_u = P_s$ (\cref{thm:b3:modules:frobenius}); equality of
degrees forces $s = 1$.

(iii)$\Rightarrow$(i): $V \cong K[X]/(P_1)$ is cyclic, generated
by the preimage of $\bar 1$.

A companion matrix is the case $V = K[X]/(P)$ itself: $x = \bar
1$, i.e.\ $e_1$, is cyclic. For a diagonal matrix
$\operatorname{diag}(\lambda_1, \dots, \lambda_n)$: $\chi = \prod
(X - \lambda_i)$, $\mu = \prod_{\lambda \text{ distinct}} (X -
\lambda)$; they agree iff the $\lambda_i$ are pairwise distinct:
a diagonal matrix has a cyclic vector iff its diagonal entries
are pairwise distinct (then $x = (1, \dots, 1)$ works:
Vandermonde).
\end{solution}

\begin{solution}{exo:b3:modules:10}
Smith: $M = Q\operatorname{diag}(d_1,\dots,d_n)P$;
\cref{exo:b3:modules:4} gives $[\Z^n : M\Z^n] = \prod\abs{d_i} =
\abs{\det M}$. If $\det M = \pm 1$: the adjugate formula
$M^{-1} = (\det M)^{-1}\,{}^{t}\!\operatorname{com}(M)$ has
integer entries, so $M \in GL_n(\Z)$; conversely $MM^{-1} = I$
gives $\det M \cdot \det M^{-1} = 1$ in $\Z$, so $\det M = \pm1$.

Subgroups of index $m$ in $\Z^2$: such a subgroup $L$ has rank
$2$ (finite index) and a unique basis in \emph{Hermite normal
form} $\bigl(\begin{smallmatrix} a & b\\ 0 &
d\end{smallmatrix}\bigr)$: $d$ is characterized by $L \cap (\{0\}
\times \Z) = \{0\} \times d\Z$, $a$ by $\pi_1(L) = a\Z$ (first
coordinates), and $b$ is then unique modulo $d$; normalize $a, d
> 0$ and $0 \leq b < d$. The index is $ad = m$. Counting: for
each divisor $d \mid m$ ($a = m/d$), there are $d$ choices of
$b$: total $\sum_{d \mid m} d = \sigma_1(m)$.
\end{solution}

\begin{solution}{pb:b3:modules:1}
\textbf{1.} $\varphi$ is additive, and $K[X]$-linear: $\varphi(X
\cdot (Q_i)_i) = \sum_i (XQ_i)(M)e_i = M\sum_i Q_i(M)e_i = X
\cdot \varphi\bigl((Q_i)_i\bigr)$, the module structure of $V_M$
being $X \cdot v = Mv$. It is surjective: constant vectors give
all of $K^n$. Column $j$ of $XI - M$ is $Xe_j - \sum_i
m_{ij}e_i$, whose image is $Me_j - \sum_i m_{ij}e_i = 0$.

\textbf{2.} Modulo the columns, $Xe_j \equiv \sum_i m_{ij}e_i$:
any vector of polynomials reduces, by induction on the top
degree, to a \emph{constant} vector $c \in K^n$. If the original
vector is in $\ker\varphi$, then $\varphi(c) = c = 0$ (on
constants, $\varphi$ is the identification $K^n = V_M$), so the
vector lies in the column span: $\ker\varphi = (XI_n -
M)K[X]^n$. Hence $V_M \cong K[X]^n/(XI - M)K[X]^n$, and Smith
over $K[X]$ (all invariant factors nonzero, their product being
$\det(XI - M) = \chi_M$) gives $V_M \cong \bigoplus_i
K[X]/(f_i)$: the nonconstant $f_i$ are the similarity
invariants, computable as quotients of minor gcds
(\cref{thm:b3:modules:smith}).

\textbf{3.} $\lambda I_n$: $XI - \lambda I$ is already Smith:
invariants $(X - \lambda, \dots, X - \lambda)$, $n$ of them.
Distinct diagonal entries: $V \cong \bigoplus_i K[X]/(X -
\lambda_i)$ with pairwise comaximal moduli, so CRT compresses to
the single cyclic $K[X]/\bigl(\prod_i(X - \lambda_i)\bigr)$:
one invariant, $\chi$. $J_n(0)$: $\mu = X^n = \chi$ forces a
single invariant $X^n$. $\operatorname{diag}(J_2(0), J_1(0))$:
elementary divisors $X^2, X$: invariants $P_1 = X \mid P_2 =
X^2$.

\textbf{4.} $P \mapsto P(u)$ maps $K[X]$ onto $K[u]$, with
kernel $(\mu_u)$ by definition of the minimal polynomial: $K[u]
\cong K[X]/(\mu_u)$, of dimension $\deg\mu_u = \deg P_s = n_s$.

\textbf{5.} $K[X]$-linearity forces $f(\bar Q) = f(Q\cdot\bar 1)
= Q\,f(\bar 1)$. The class $\bar 1$ satisfies $P \bar 1 = 0$, so
$P f(\bar 1) = 0$ is necessary. Conversely if $P\bar c = 0$,
then $f(\bar Q) = Q\bar c$ is well defined ($Q \equiv Q' \bmod P
\Rightarrow (Q - Q')\bar c \in$ multiples of $P\bar c = 0$) and
$K[X]$-linear.

\textbf{6.} Let $g = \gcd(P, Q)$, $P = gP'$, $Q = gQ'$ with
$\gcd(P', Q') = 1$. In $K[X]/(Q)$: $P\bar c = 0 \iff Q \mid Pc
\iff Q' \mid P'c \iff Q' \mid c$ (Euclid, $\gcd(P', Q') = 1$).
So the admissible $\bar c$ form the submodule generated by
$\bar{Q'}$, whose annihilator is $\{R : Q \mid RQ'\} = (g)$:
that submodule is $\cong K[X]/(g)$. With question 5,
$\operatorname{Hom}(K[X]/(P), K[X]/(Q)) \cong K[X]/(\gcd(P,Q))$,
of dimension $\deg\gcd(P,Q)$.

\textbf{7.} $v$ commutes with $u$ iff $v$ commutes with every
$P(u)$, iff $v$ is $K[X]$-linear: $\mathcal C(u) =
\operatorname{End}_{K[X]}(V)$. Writing morphisms of $V =
\bigoplus_j V_j$ as matrices $(f_{ij})$, $f_{ij} \in
\operatorname{Hom}(V_j, V_i)$ (compose with injections and
projections), question 6 gives
\[
\dim \mathcal C(u) = \sum_{i,j} \deg\gcd(P_i, P_j)
= \sum_{i,j} n_{\min(i,j)}
= \sum_{k=1}^{s} \bigl(2(s - k) + 1\bigr)\,n_k ,
\]
using the divisibility chain ($\gcd(P_i, P_j) =
P_{\min(i,j)}$) and, for the last step, that $\min(i,j) = k$
happens for exactly $2(s - k) + 1$ pairs $(i, j)$.

\textbf{8.} Since $2(s-k)+1 \geq 1$, $\dim\mathcal C(u) \geq
\sum_k n_k = n$, with equality iff $s = 1$, i.e.\ iff $u$ is
cyclic (\cref{exo:b3:modules:9}). For $u = \lambda\,\mathrm{id}$:
$s = n$, all $n_k = 1$: $\dim = \sum_{k=1}^n (2(n-k)+1) = n^2$
--- correct, since $\mathcal C(\lambda\,\mathrm{id}) = \mathcal
L(V)$.

\textbf{9.} Formula: invariants $(X, X^2)$, so $s = 2$, $n_1 =
1$, $n_2 = 2$: $\dim = 3\cdot1 + 1\cdot2 = 5$. Directly: in the
basis $(e_1, e_2, e_3)$ with $u e_2 = e_1$, $ue_1 = ue_3 = 0$,
writing $Au = uA$ for $A = (a_{ij})$ yields the conditions
$a_{21} = a_{23} = a_{31} = 0$ and $a_{11} = a_{22}$: five free
parameters $a_{11}{=}a_{22}, a_{12}, a_{13}, a_{32}, a_{33}$.

\textbf{10.} A polynomial $P(u)$ commutes with anything that
commutes with $u$ (it is a sum of powers of $u$): $K[u]
\subseteq \mathcal C(\mathcal C(u))$. And $u \in \mathcal C(u)$,
so any $w \in \mathcal C(\mathcal C(u))$ commutes with $u$.

\textbf{11.} Let $V = K[u]x$ and $v \in \mathcal C(u)$. Write
$v(x) = P(u)x$ (cyclicity). For $y = Q(u)x$ arbitrary: $v(y) =
vQ(u)x = Q(u)v(x) = Q(u)P(u)x = P(u)y$. So $v = P(u)$: $\mathcal
C(u) = K[u]$. Then $\mathcal C(\mathcal C(u)) = \mathcal
C(K[u]) = \mathcal C(u) = K[u]$ (commuting with all of $K[u]$ is
the same as commuting with $u$). The theorem holds in the cyclic
case.

\textbf{12.} $\pi_i$ is $K[X]$-linear (the decomposition is a
direct sum of submodules), so $\pi_i \in \mathcal C(u)$, and $w$
commutes with it: $w(V_i) = w\pi_i(V) = \pi_i w(V) \subseteq
V_i$. The restriction $w_i = w\restriction_{V_i}$ commutes with
the cyclic $u_i = u\restriction_{V_i}$ (question 10), and $w_i
\in \mathcal C(u_i) = K[u_i]$ (question 11): $w_i = Q_i(u_i) =
Q_i(u)\restriction_{V_i}$ for some $Q_i \in K[X]$.

\textbf{13.} Well-definedness of $\eta_{ij}(R(u)x_j) = R(u)x_i$:
if $R(u)x_j = 0$ then $P_j \mid R$, and $P_i \mid P_j$ gives
$P_i \mid R$, so $R(u)x_i = 0$ ($\operatorname{Ann}(x_i) =
(P_i)$). $\eta_{ij}$ is then $K[X]$-linear by construction, and
$\tilde\eta_{ij} = \eta_{ij}\pi_j$ is a composition of
$K[X]$-morphisms $V \to V$: $\tilde\eta_{ij} \in \mathcal
C(u)$.

\textbf{14.} Evaluate $w\tilde\eta_{ij} = \tilde\eta_{ij}w$ at
$x_j$: the left side is $w(x_i) = Q_i(u)x_i$; the right side is
$\eta_{ij}\bigl(Q_j(u)x_j\bigr) = Q_j(u)x_i$. Hence $(Q_i -
Q_j)(u)\,x_i = 0$: $P_i \mid Q_i - Q_j$ for all $i \leq j$. In
particular, with $Q = Q_s$: $Q \equiv Q_i \pmod{P_i}$, so $Q(u)$
and $Q_i(u)$ agree on $V_i$ (which $P_i(u)$ kills). Therefore $w
= Q(u)$ on each $V_i$, hence on $V$: $\mathcal C(\mathcal C(u))
\subseteq K[u]$, and with question 10, $\mathcal C(\mathcal
C(u)) = K[u]$.

\textbf{15.} The first assertion is questions 10--14 (or, for
cyclic $u$, question 11 alone). For $u = \mathrm{id}$, $n \geq
2$: $\mathcal C(u) = \mathcal L(V)$ has dimension $n^2$, while
$\dim K[u] = \deg\mu_u = 1$. So $\mathcal C(u) = K[u]$ fails
badly; yet the bicommutant theorem
holds ($\mathcal C(\mathcal L(V)) = K\,\mathrm{id} = K[u]$: the
center of the matrix algebra is the scalars). Cyclicity is what
makes the \emph{single} commutant already polynomial; the
\emph{double} commutant is polynomial always.
\end{solution}
